\section{Conclusion}
\label{sec:conclusion}

This paper presents ASCENSION, a novel DA method built on the
intuition that gradually exploring under-represented latent
regions and expanding the initial data distribution can enhance
DA effectiveness. ASCENSION leverages a VAE with a contrastive
clustering loss to promote intra-class compactness and
inter-class separability, and introduces an $\alpha$-scaling
mechanism that progressively expands class densities to generate
diverse, class-consistent samples.
% REVISED FOR TMLR (rebuttal C2/C3)
Across 102 UCR benchmarks and three classifiers (ResNet, FCN,
Moment) under a single fixed configuration, ASCENSION is the only
method with a positive mean accuracy gain on each classifier
($+1.3\%$ on ResNet, $+0.6\%$ on FCN, and $+0.8\%$ on Moment),
with non-negative gains on $71.9\%$, $62.7\%$, and $62.4\%$ of
datasets respectively, whereas the strongest baseline (VaDE) loses
$0.9\%$ on average. ASCENSION is never significantly worse than
any competitor, and its benefit concentrates on the datasets where
training data is scarcest: below 20 training samples per class the
median gain is clearly positive on ResNet and Moment, while above
50 samples per class ASCENSION is close to neutral. An ablation study confirms that the $\alpha$-scaling
mechanism is the single most impactful component of the method.
These results position ASCENSION as a low-risk augmentation default
whose target regime is the small-training-set setting, a regime
that remains underserved by generative DA, without requiring prior
knowledge of method suitability.
